\section{Preliminaries for the obstruction proofs}\label{app:obstruction-preliminaries}
The constructions below use three elementary facts about valued fields and power series. We record their proofs to make explicit the scope of the coefficient argument, including incomplete base fields.

\begin{lemma}[one-variable uniqueness]\label{lem:unique}
Let $K$ be nontrivially normed and $Z$ a normed $K$\nobreakdash-space, not necessarily complete. Let $z_0,z_1,\ldots\in Z$, $\rho>0$
and $M\ge0$ with $\norm{z_n}\rho^n\le M$ for all $n$. If $\sum_nt^nz_n$ converges to $0$ for every $t\in K$ with
$\abs t<\rho$, then $z_n=0$ for all $n$.
\end{lemma}

\begin{proof}
Otherwise let $N$ be least with $z_N\neq0$. For $0<\abs t<\rho/2$ the series $\sum_{n>N}t^{n-N}z_n$ converges, being
$t^{-N}$ times the convergent tail $\sum_{n>N}t^nz_n$, and $t^N\bigl(z_N+\sum_{n>N}t^{n-N}z_n\bigr)=0$. Hence
\[
\norm{z_N}=\Bigl\lVert\sum_{n>N}t^{n-N}z_n\Bigr\rVert\le\sum_{n>N}M\rho^{-n}\abs t^{n-N}
=M\rho^{-N}\frac{\abs t/\rho}{1-\abs t/\rho},
\]
where the norm of a convergent series is bounded by the sum of the norms by continuity of the norm on partial sums.
Since $K$ is nontrivially normed, there are $t\ne0$ with $\abs t$ arbitrarily small, so $\norm{z_N}=0$, a
contradiction.
\end{proof}

\begin{lemma}\label{lem:ultra}
A normed field of characteristic $p>0$ is nonarchimedean.
\end{lemma}

\begin{proof}
Frobenius and the triangle inequality give, for every $m\ge0$,
\[
 \abs{x+y}^{p^m}=\abs{x^{p^m}+y^{p^m}}
 \le2\max(\abs x,\abs y)^{p^m}.
\]
Taking $p^m$\nobreakdash-th roots and letting $m\to\infty$ proves the assertion.
\end{proof}

\begin{lemma}\label{lem:findim}
Let $\mathbb K$ be a complete nontrivially normed field and $U$ a finite-dimensional normed $\mathbb K$\nobreakdash-space. Then every
linear functional on $U$ is bounded, and $U$ is complete.
\end{lemma}

\begin{proof}
We prove both assertions by induction on $n=\dim U$, the case $n=0$ being immediate. Suppose $n\ge1$ and let $g\ne0$
be a linear functional on $U$. Its kernel $H$ has dimension $n-1$, so the induction hypothesis makes $H$ complete
and therefore closed in $U$. Choose $v$ with $g(v)=1$. Since $v\notin H$,
$\delta:=\inf_{h\in H}\norm{v-h}>0$. Write $x=g(x)v+h$ with $h\in H$. If $g(x)\ne0$, then
\[
\norm x=\abs{g(x)}\,\norm{v+g(x)^{-1}h}\ge\abs{g(x)}\delta.
\]
Thus $\abs{g(x)}\le\delta^{-1}\norm x$ for every $x$, proving boundedness.

Choose a basis of $U$. Its coordinate functionals are now bounded, so the coordinate map
$U\to\mathbb K^n$, with the max norm on $\mathbb K^n$, is bounded. The inverse map
$(c_i)\mapsto\sum_ic_ie_i$ is bounded by $\sum_i\norm{e_i}$. Hence $U$ is bi-Lipschitz equivalent to the complete
space $\mathbb K^n$, and is itself complete.
\end{proof}

When the scalar field needs to be displayed, $A^{k,L}$ denotes the degree-$k$ precomposition action formed over $L$. For the precomposition action $A$, a \emph{bounded lift} means a bounded $k$\nobreakdash-linear map on its operator domain whose diagonal is $A$, as in Proposition~\ref{prop:homogeneous-lift}. That proposition applies at every base point; Lemma~\ref{lem:unique} provides the one-variable coefficient comparison used in its proof.
